\documentclass[aspectratio=169,10pt]{beamer}

\usetheme{Madrid}
\usefonttheme{professionalfonts}
\usepackage[T1]{fontenc}
\usepackage[utf8]{inputenc}
\usepackage{lmodern}
\usepackage{listings}
\usepackage{xcolor}
\usepackage{booktabs}
\usepackage{array}
\usepackage{hyperref}
\usepackage{ragged2e}

% Palette
\definecolor{Navy}{HTML}{17324D}
\definecolor{Blue}{HTML}{2D6A9F}
\definecolor{PaleBlue}{HTML}{EAF2F8}
\definecolor{Coral}{HTML}{E76F68}
\definecolor{PaleCoral}{HTML}{FCEDEC}
\definecolor{Ink}{HTML}{24313D}
\definecolor{Muted}{HTML}{5E6B75}
\definecolor{CodeBG}{HTML}{F4F7FA}
\definecolor{CodeGreen}{HTML}{26734D}
\definecolor{CodePurple}{HTML}{7652A3}

\setbeamercolor{structure}{fg=Navy}
\setbeamercolor{palette primary}{bg=Navy,fg=white}
\setbeamercolor{palette secondary}{bg=Blue,fg=white}
\setbeamercolor{palette tertiary}{bg=Navy!88!black,fg=white}
\setbeamercolor{palette quaternary}{bg=Coral,fg=white}
\setbeamercolor{title}{fg=Navy}
\setbeamercolor{frametitle}{fg=Navy,bg=PaleBlue}
\setbeamercolor{block title}{fg=white,bg=Navy}
\setbeamercolor{block body}{fg=Ink,bg=PaleBlue}
\setbeamercolor{alerted text}{fg=Coral}
\setbeamercolor{item projected}{bg=Coral,fg=white}
\setbeamercolor{normal text}{fg=Ink,bg=white}
\setbeamercolor{caption name}{fg=Navy}

\setbeamertemplate{navigation symbols}{}
\setbeamertemplate{footline}{
  \leavevmode%
  \hbox{%
  \begin{beamercolorbox}[wd=.78\paperwidth,ht=2.6ex,dp=1.1ex,leftskip=1em]{author in head/foot}%
    \color{Muted}\insertshorttitle
  \end{beamercolorbox}%
  \begin{beamercolorbox}[wd=.22\paperwidth,ht=2.6ex,dp=1.1ex,rightskip=1em plus1fil]{date in head/foot}%
    \color{Muted}\insertframenumber{} / \inserttotalframenumber
  \end{beamercolorbox}}%
  \vskip0pt%
}
\setbeamertemplate{frametitle}{
  \nointerlineskip
  \begin{beamercolorbox}[wd=\paperwidth,ht=3.1ex,dp=1.1ex,leftskip=0.8em,rightskip=0.5em]{frametitle}
    \usebeamerfont{frametitle}\insertframetitle
  \end{beamercolorbox}
  \vspace{0.15em}
}
\setbeamersize{text margin left=7mm,text margin right=7mm}

\lstdefinestyle{python}{
  language=Python,
  backgroundcolor=\color{CodeBG},
  basicstyle=\ttfamily\scriptsize\color{Ink},
  keywordstyle=\color{Blue}\bfseries,
  stringstyle=\color{CodeGreen},
  commentstyle=\color{Muted}\itshape,
  numberstyle=\tiny\color{Muted},
  numbers=none,
  showstringspaces=false,
  breaklines=true,
  breakatwhitespace=true,
  frame=single,
  rulecolor=\color{PaleBlue},
  frameround=tttt,
  tabsize=4,
  keepspaces=true,
  columns=fullflexible,
  xleftmargin=3pt,
  xrightmargin=3pt,
  aboveskip=4pt,
  belowskip=4pt
}
\lstset{style=python}

\newcommand{\scenario}[1]{\begin{block}{Real-world scenario}#1\end{block}}
\newcommand{\whenuse}[1]{\textcolor{CodeGreen}{\textbf{When to use:}} #1}
\newcommand{\pitfall}[1]{\begin{alertblock}{Watch out}#1\end{alertblock}}
\newcommand{\method}[1]{\texttt{\textcolor{Blue}{#1}}}

\title[Pandas: Beginner to Advanced]{Python pandas\\Beginner to Advanced}
\subtitle{A practical, scenario-based guide to cleaning, exploring, transforming, and analyzing data}
\author{Nafiul Islam}
\institute{BSDS \\ United International University}
\date{\today}

\begin{document}

% 1
\begin{frame}[plain,fragile]
  \titlepage
  \vspace{-0.6em}
  \begin{center}
    \small\textcolor{Muted}{Learn the tools. Understand the choices. Solve real data problems.}
  \end{center}
\end{frame}

% 2
\begin{frame}[fragile]{How to use this course}
\begin{columns}[T,onlytextwidth]
\column{0.49\textwidth}
\begin{block}{Learning path}
\begin{enumerate}
\item Understand tables and pandas objects
\item Load, inspect, and select data
\item Clean unreliable data
\item Analyze and summarize patterns
\item Combine and reshape datasets
\item Work with dates, text, and categories
\item Build reliable, efficient workflows
\end{enumerate}
\end{block}
\column{0.49\textwidth}
\begin{block}{How each example works}
\begin{itemize}
\item \textbf{Scenario:} a real problem
\item \textbf{Code:} a small working pattern
\item \textbf{Meaning:} what the code does
\item \textbf{Choice:} when to use it
\item \textbf{Caution:} common mistakes
\end{itemize}
\end{block}
\end{columns}
\small This is a broad learning deck, not a replacement for the API reference. Use the official documentation links near the end when you need every parameter or edge case.
\end{frame}

% 3
\begin{frame}[fragile]{What is pandas?}
\begin{itemize}
\item \textbf{pandas} is a Python library for working with structured data: tables, columns, dates, and labels.
\item It is useful when data comes from spreadsheets, CSV files, databases, APIs, or experiments.
\item It helps you inspect, clean, filter, combine, summarize, and reshape data.
\item It works especially well with NumPy and can plot through Matplotlib.
\end{itemize}
\scenario{A school has 2,000 student records. You want to find average marks, missing phone numbers, and the highest-scoring class without manually editing a spreadsheet.}
\begin{lstlisting}
import pandas as pd
\end{lstlisting}
\whenuse{Use pandas when your problem involves rows and columns. It is not a replacement for every database or distributed big-data system.}
\end{frame}

% 4
\begin{frame}[fragile]{Install pandas and start a project}
\begin{lstlisting}
# In a terminal
python -m pip install pandas

# Optional tools for charts and notebooks
python -m pip install matplotlib jupyter openpyxl

# In Python
import pandas as pd
print(pd.__version__)
\end{lstlisting}
\begin{block}{Why use \texttt{pd}?}
\texttt{pd} is the common short alias for pandas. It keeps code shorter and is widely recognized by Python developers.
\end{block}
\pitfall{Install packages in the same Python environment that runs your notebook or script. If \texttt{import pandas} fails, check the active environment first.}
\end{frame}

% 5
\begin{frame}[fragile]{The two main pandas data structures}
\begin{columns}[T,onlytextwidth]
\column{0.48\textwidth}
\begin{block}{Series}
A single labeled column.
\begin{lstlisting}
marks = pd.Series(
    [78, 91, 84],
    index=["Ava", "Rafi", "Mina"],
    name="marks"
)
\end{lstlisting}
Think: one spreadsheet column with labels.
\end{block}
\column{0.48\textwidth}
\begin{block}{DataFrame}
A table made of rows and columns.
\begin{lstlisting}
students = pd.DataFrame({
    "name": ["Ava", "Rafi", "Mina"],
    "marks": [78, 91, 84]
})
\end{lstlisting}
Think: a full spreadsheet or database table.
\end{block}
\end{columns}
\scenario{Use a Series for one measure such as daily temperature. Use a DataFrame for a weather table containing date, city, temperature, and rainfall.}
\end{frame}

% 6
\begin{frame}[fragile]{Create a small practice dataset}
\begin{lstlisting}
sales = pd.DataFrame({
    "order_id": [101, 102, 103, 104, 105, 106],
    "city": ["Dhaka", "Sylhet", "Dhaka", "Khulna", "Sylhet", "Dhaka"],
    "product": ["Pen", "Book", "Bag", "Pen", "Bag", "Book"],
    "units": [3, 2, 1, 5, 2, 4],
    "price": [20, 150, 800, 20, 800, 150],
    "discount": [0, 10, 50, 0, 0, 15]
})
sales["revenue"] = sales["units"] * sales["price"] - sales["discount"]
\end{lstlisting}
\begin{itemize}
\item Each row represents one order.
\item Each column describes one feature of the order.
\item The \texttt{revenue} column is calculated from existing columns.
\end{itemize}
\scenario{We will reuse this fictional store dataset to learn methods in a connected, realistic workflow.}
\end{frame}

% 7
\begin{frame}[fragile]{Read a table before changing it}
\begin{lstlisting}
sales.head()       # first 5 rows
sales.head(10)     # first 10 rows
sales.tail(3)      # last 3 rows
sales.sample(3)    # 3 random rows
sales.shape        # (rows, columns)
sales.columns      # column labels
sales.index        # row labels
sales.dtypes       # each column's data type
\end{lstlisting}
\begin{block}{Which one should I use?}
\begin{itemize}
\item \method{head()} / \method{tail()}: quick visual checks.
\item \method{shape}: size of the dataset.
\item \method{dtypes}: find numbers stored as text or dates stored as strings.
\item \method{sample()}: inspect varied rows rather than only the beginning.
\end{itemize}
\end{block}
\end{frame}

% 8
\begin{frame}[fragile]{Inspect quality and get a quick summary}
\begin{lstlisting}
sales.info()
sales.describe()                 # numeric columns by default
sales.describe(include="all")    # broader summary
sales["city"].value_counts()
sales["city"].nunique()
sales.isna().sum()
\end{lstlisting}
\scenario{A downloaded sales file has 80,000 rows. Before analysis, you need to know whether the price column is numeric, how many cities appear, and whether any fields are missing.}
\whenuse{Start every unfamiliar dataset with \method{head()}, \method{info()}, \method{describe()}, and missing-value counts.}
\pitfall{\method{describe()} is a summary, not proof that the data is correct. Impossible values can still have valid data types.}
\end{frame}

% 9
\begin{frame}[fragile]{Load data from common file formats}
\begin{lstlisting}
df = pd.read_csv("orders.csv")
df = pd.read_csv("orders.csv", sep=";")
df = pd.read_excel("orders.xlsx", sheet_name="Orders")
df = pd.read_json("orders.json")
df = pd.read_parquet("orders.parquet")

df.to_csv("clean_orders.csv", index=False)
df.to_excel("clean_orders.xlsx", index=False)
df.to_json("clean_orders.json", orient="records")
df.to_parquet("clean_orders.parquet", index=False)
\end{lstlisting}
\begin{itemize}
\item \method{read\_...()} loads data; \method{to\_...()} writes data.
\item CSV is simple and widely supported; Parquet is often more efficient for analytical workflows.
\item Excel support may require an optional engine such as \texttt{openpyxl}.
\end{itemize}
\pitfall{Do not overwrite the only copy of raw data. Keep a raw file and save cleaned output separately.}
\end{frame}

% 10
\begin{frame}[fragile]{Control how data is imported}
\begin{lstlisting}
df = pd.read_csv(
    "orders.csv",
    usecols=["order_id", "city", "price"],
    nrows=1000,
    dtype={"order_id": "string"},
    na_values=["", "N/A", "unknown"]
)
\end{lstlisting}
\begin{itemize}
\item \method{usecols}: load only needed columns.
\item \method{nrows}: read a small sample while testing.
\item \method{dtype}: request a column type, when appropriate.
\item \method{na\_values}: tell pandas which tokens mean missing data.
\end{itemize}
\scenario{A large transaction file contains 60 columns, but your first task only needs order ID, city, and price. Loading fewer columns can save memory and time.}
\end{frame}

% 11
\begin{frame}[fragile]{Series and DataFrame: useful properties}
\begin{lstlisting}
sales.shape
sales.size
sales.ndim
sales.empty
sales.memory_usage(deep=True)
sales["price"].name
sales["price"].dtype
sales["price"].index
\end{lstlisting}
\begin{block}{Remember}
\begin{itemize}
\item \method{shape} returns rows and columns.
\item \method{size} counts cells.
\item \method{ndim} gives number of dimensions.
\item \method{memory\_usage()} helps investigate memory use.
\end{itemize}
\end{block}
\whenuse{Use these properties for quick diagnostics and for checking that a transformation produced the expected table size.}
\end{frame}

% 12
\begin{frame}[fragile]{Select one column or several columns}
\begin{lstlisting}
sales["city"]                         # Series
sales[["city", "product"]]            # DataFrame
sales.loc[:, ["order_id", "revenue"]] # all rows, chosen columns
\end{lstlisting}
\begin{itemize}
\item One pair of brackets with a column name normally returns a Series.
\item A list of column names returns a DataFrame.
\item Use explicit brackets rather than relying on \texttt{sales.city}; attribute access can conflict with method names or unusual column labels.
\end{itemize}
\scenario{A manager only needs city and revenue for a quick report. Select those columns rather than carrying every field into the next step.}
\end{frame}

% 13
\begin{frame}[fragile]{Select rows with \texttt{loc} and \texttt{iloc}}
\begin{columns}[T,onlytextwidth]
\column{0.49\textwidth}
\begin{block}{\texttt{loc}: by label}
\begin{lstlisting}
sales.loc[0]
sales.loc[0:2, ["city", "price"]]
\end{lstlisting}
Use labels from the index. Label slices include the end label when present.
\end{block}
\column{0.49\textwidth}
\begin{block}{\texttt{iloc}: by position}
\begin{lstlisting}
sales.iloc[0]
sales.iloc[0:3, 1:4]
\end{lstlisting}
Use integer positions. Standard Python-style slices exclude the stop position.
\end{block}
\end{columns}
\pitfall{The index label is not always the row's current position. After sorting or filtering, label 10 may not be the eleventh row.}
\end{frame}

% 14
\begin{frame}[fragile]{Filter rows with conditions}
\begin{lstlisting}
sales[sales["city"] == "Dhaka"]
sales[sales["price"] > 100]
sales[(sales["city"] == "Dhaka") & (sales["units"] >= 2)]
sales[(sales["product"] == "Bag") | (sales["price"] < 100)]
sales[~(sales["city"] == "Sylhet")]
\end{lstlisting}
\begin{itemize}
\item \texttt{\&} means AND; \texttt{|} means OR; \texttt{\~{}} means NOT.
\item Put each condition in parentheses.
\item Use \texttt{==} to compare; a single \texttt{=} assigns a value in Python.
\end{itemize}
\scenario{A shop manager wants all Dhaka orders with at least two units to plan local delivery.}
\end{frame}

% 15
\begin{frame}[fragile]{More useful selection patterns}
\begin{lstlisting}
sales[sales["city"].isin(["Dhaka", "Sylhet"])]
sales[sales["price"].between(100, 500)]
sales.query("units >= 2 and price < 200")
sales.nlargest(3, "revenue")
sales.nsmallest(2, "price")
sales.sample(n=2, random_state=7)
\end{lstlisting}
\begin{itemize}
\item \method{isin}: match any item from a list.
\item \method{between}: test a numeric range.
\item \method{query}: express readable conditions as a string.
\item \method{nlargest}/\method{nsmallest}: retrieve top or bottom rows by a column.
\end{itemize}
\scenario{Find low-cost products for a promotion, or the three orders with the highest revenue.}
\pitfall{\method{query()} is convenient for simple expressions, but normal boolean indexing is more flexible for complex logic.}
\end{frame}

% 16
\begin{frame}[fragile]{Sort rows and control the index}
\begin{lstlisting}
sales.sort_values("revenue")
sales.sort_values("revenue", ascending=False)
sales.sort_values(["city", "revenue"], ascending=[True, False])
sales.sort_index()
sales.reset_index(drop=True)
sales.set_index("order_id")
\end{lstlisting}
\begin{itemize}
\item \method{sort\_values}: sort by column values.
\item \method{sort\_index}: sort by row or column labels.
\item \method{set\_index}: make a column the row label.
\item \method{reset\_index}: return to a simple numbered index.
\end{itemize}
\scenario{Rank each city's orders by revenue, or use order ID as a meaningful lookup label.}
\end{frame}

% 17
\begin{frame}[fragile]{Create, change, and remove columns}
\begin{lstlisting}
sales["gross"] = sales["units"] * sales["price"]
sales["is_large_order"] = sales["units"] >= 4
sales = sales.assign(
    average_unit_price=sales["revenue"] / sales["units"]
)
sales = sales.rename(columns={"city": "customer_city"})
sales = sales.drop(columns=["discount"])
\end{lstlisting}
\begin{itemize}
\item Assignment creates or replaces a column.
\item \method{assign()} can make a readable chain of transformations.
\item \method{rename()} changes labels.
\item \method{drop()} removes rows or columns; state \texttt{columns=} to be explicit.
\end{itemize}
\pitfall{Check your business definition before calculating a metric. In real sales data, tax, shipping, refunds, and discount rules may affect revenue.}
\end{frame}

% 18
\begin{frame}[fragile]{Vectorized calculations: avoid unnecessary loops}
\begin{lstlisting}
# Recommended: calculate an entire column at once
sales["gross"] = sales["units"] * sales["price"]
sales["discount_rate"] = sales["discount"] / (
    sales["units"] * sales["price"]
)

# Usually slower and harder to maintain:
# loop through every row and calculate one by one
\end{lstlisting}
\begin{block}{What does vectorized mean?}
You describe an operation on a whole Series or column. pandas and its underlying libraries handle the element-by-element work efficiently.
\end{block}
\scenario{Calculate total cost for 500,000 shopping-cart rows without writing a Python loop for every row.}
\pitfall{Vectorization does not mean every expression is magically fast. Avoid creating many huge temporary columns and benchmark important workloads.}
\end{frame}

% 19
\begin{frame}[fragile]{Missing values: find them first}
\begin{lstlisting}
sales.isna()
sales.isna().sum()
sales.notna()
sales["price"].isna()
sales.dropna()
sales.dropna(subset=["price"])
\end{lstlisting}
\begin{itemize}
\item Missing values may appear as \texttt{NaN}, \texttt{None}, or pandas' \texttt{NA}, depending on type and source.
\item \method{isna()} marks missing cells; \method{notna()} marks present cells.
\item \method{dropna()} removes rows or columns according to its settings.
\end{itemize}
\scenario{A price is missing from 12 product records. Find those records before deciding whether to correct, remove, or keep them.}
\pitfall{Missing does not always mean zero. A missing age, a missing price, and a missing survey answer have different meanings.}
\end{frame}

% 20
\begin{frame}[fragile]{Fill missing values thoughtfully}
\begin{lstlisting}
sales["discount"] = sales["discount"].fillna(0)
sales["price"] = sales["price"].fillna(
    sales["price"].median()
)
sales["city"] = sales["city"].fillna("Unknown")

# Forward/backward fill: useful for ordered measurements
weather["temperature"] = weather["temperature"].ffill()
weather["temperature"] = weather["temperature"].bfill()
\end{lstlisting}
\begin{itemize}
\item \method{fillna(0)} only when zero is the real meaning.
\item Median is less affected by extreme values than mean.
\item \method{ffill()} carries the previous valid value forward; \method{bfill()} uses the next valid value.
\end{itemize}
\pitfall{For time series, sort by time first. Never fill using future information if you are simulating a real-time prediction.}
\end{frame}

% 21
\begin{frame}[fragile]{Data types and safe conversion}
\begin{lstlisting}
df.dtypes
df["price"] = pd.to_numeric(df["price"], errors="coerce")
df["order_date"] = pd.to_datetime(df["order_date"], errors="coerce")
df["customer_id"] = df["customer_id"].astype("string")
df["units"] = df["units"].astype("Int64")
\end{lstlisting}
\begin{itemize}
\item \method{to\_numeric}: parse numeric values; \texttt{coerce} turns invalid entries into missing values.
\item \method{to\_datetime}: parse dates and times.
\item \method{astype}: convert to a requested dtype when the data supports it.
\item Nullable \texttt{Int64} can store integers and missing values together.
\end{itemize}
\scenario{A CSV has prices like \texttt{"250"} and \texttt{"unknown"}. Convert to numbers, then investigate the rows that became missing.}
\end{frame}

% 22
\begin{frame}[fragile]{Text cleaning with the \texttt{str} accessor}
\begin{lstlisting}
customers["name"] = customers["name"].str.strip()
customers["city"] = customers["city"].str.title()
customers["email"] = customers["email"].str.lower()
customers["phone"] = customers["phone"].str.replace(" ", "", regex=False)

customers["has_gmail"] = customers["email"].str.contains(
    "@gmail.com", na=False, regex=False
)
customers["domain"] = customers["email"].str.split("@").str[-1]
\end{lstlisting}
\begin{itemize}
\item \method{str.strip}: remove leading/trailing spaces.
\item \method{str.lower}/\method{str.upper}/\method{str.title}: normalize letter case.
\item \method{str.replace}: replace text.
\item \method{str.contains}: test whether a pattern occurs.
\item \method{str.split}: split text into pieces.
\end{itemize}
\scenario{The same city arrives as \texttt{"dhaka"}, \texttt{"Dhaka "}, and \texttt{"DHAKA"}. Normalize it before counting customers by city.}
\end{frame}

% 23
\begin{frame}[fragile]{Extract patterns and validate text}
\begin{lstlisting}
orders["order_code"].str.startswith("BD-")
orders["email"].str.endswith(".com", na=False)
orders["code_digits"] = orders["order_code"].str.extract(r"(\d+)")
orders["clean_code"] = orders["order_code"].str.replace(
    r"[^A-Za-z0-9-]", "", regex=True
)
\end{lstlisting}
\begin{itemize}
\item \method{startswith}/\method{endswith}: check the beginning or end.
\item \method{extract}: capture part of a string with a regular expression.
\item Regex is a compact language for describing text patterns.
\end{itemize}
\scenario{An inventory code such as \texttt{ITEM-00482} contains a numeric product identifier that must be extracted for matching.}
\pitfall{A pattern match is not the same as full validation. For emails, IDs, or phone numbers, define the actual rules your organization requires.}
\end{frame}

% 24
\begin{frame}[fragile]{Replace values and standardize categories}
\begin{lstlisting}
df["status"] = df["status"].replace({
    "complete": "Completed",
    "done": "Completed",
    "pending ": "Pending"
})
df["rating"] = df["rating"].replace({-1: pd.NA})
df["status"] = df["status"].str.strip().str.title()
\end{lstlisting}
\begin{itemize}
\item \method{replace()} changes exact values or matching patterns.
\item \method{str.strip()} removes accidental spaces.
\item Standardize spelling before grouping or joining.
\end{itemize}
\scenario{A delivery dataset uses \texttt{"done"}, \texttt{"complete"}, and \texttt{"Completed"} for the same outcome. Standardize these labels before calculating delivery rates.}
\end{frame}

% 25
\begin{frame}[fragile]{Duplicates: detect before deleting}
\begin{lstlisting}
df.duplicated()
df.duplicated().sum()
df[df.duplicated(subset=["email"], keep=False)]
df = df.drop_duplicates()
df = df.drop_duplicates(subset=["email"], keep="last")
\end{lstlisting}
\begin{itemize}
\item \method{duplicated()} marks rows that repeat earlier rows under the selected rule.
\item \method{drop\_duplicates()} keeps one row from each duplicate group.
\item \texttt{keep="first"}, \texttt{"last"}, or \texttt{False} control which rows remain.
\end{itemize}
\scenario{A customer import contains the same email twice. Inspect timestamps and source records before deciding which version to keep.}
\pitfall{Two identical purchases may be real separate purchases. Do not delete duplicates without understanding what makes a record unique.}
\end{frame}

% 26
\begin{frame}[fragile]{Use conditions to label or replace values}
\begin{lstlisting}
sales["order_size"] = "Small"
sales.loc[sales["units"] >= 3, "order_size"] = "Large"

sales["priority"] = sales["revenue"].where(
    sales["revenue"] >= 300, other=0
)
\end{lstlisting}
\begin{itemize}
\item \method{loc} is a clear way to set values for rows meeting a condition.
\item \method{where(condition, other)} keeps values where the condition is true and substitutes elsewhere.
\item For several categories, consider \method{np.select()} or a small, well-tested function.
\end{itemize}
\scenario{A warehouse flags orders with at least three units as large so staff can choose a suitable package.}
\pitfall{Avoid chained assignment such as \texttt{df[df["x"] > 0]["y"] = 1}. Select rows and columns together with \texttt{loc}.}
\end{frame}

% 27
\begin{frame}[fragile]{Map values and create categories}
\begin{lstlisting}
city_region = {
    "Dhaka": "Central",
    "Sylhet": "Northeast",
    "Khulna": "Southwest"
}
sales["region"] = sales["city"].map(city_region)

sales["price_band"] = pd.cut(
    sales["price"],
    bins=[0, 50, 200, 1000],
    labels=["Low", "Medium", "High"],
    include_lowest=True
)
\end{lstlisting}
\begin{itemize}
\item \method{map}: look up a replacement or related value using a dictionary or function.
\item \method{cut}: place numeric values into ranges.
\item Values not present in the mapping may become missing.
\end{itemize}
\scenario{Turn a product price into Low, Medium, or High bands for a simple sales dashboard.}
\end{frame}

% 28
\begin{frame}[fragile]{Basic statistics: understand the numbers}
\begin{lstlisting}
sales["revenue"].sum()
sales["revenue"].mean()
sales["revenue"].median()
sales["revenue"].min()
sales["revenue"].max()
sales["revenue"].std()
sales["revenue"].quantile([0.25, 0.50, 0.75])
sales["city"].value_counts(normalize=True)
\end{lstlisting}
\begin{itemize}
\item \method{mean}: arithmetic average; can be pulled by extreme values.
\item \method{median}: middle value after sorting; often useful for skewed prices.
\item \method{std}: describes spread around the mean.
\item \method{quantile}: shows cut-off points such as the 25th and 75th percentiles.
\item \method{value\_counts(normalize=True)} returns proportions instead of raw counts.
\end{itemize}
\end{frame}

% 29
\begin{frame}[fragile]{Count unique values and check distributions}
\begin{lstlisting}
sales["city"].count()       # non-missing entries
sales["city"].nunique()     # distinct non-missing values
sales["city"].unique()      # distinct values as an array
sales["city"].value_counts(dropna=False)
\end{lstlisting}
\scenario{A survey contains 1,500 responses. You want the number of valid city entries, the number of different cities, and a count for each city including missing entries.}
\pitfall{\method{count()} is not the same as \method{len(df)}: count on a Series excludes missing values.}
\end{frame}

% 30
\begin{frame}[fragile]{GroupBy: split, calculate, combine}
\begin{lstlisting}
sales.groupby("city")["revenue"].sum()
sales.groupby("product")["units"].mean()
sales.groupby(["city", "product"])["revenue"].sum()
\end{lstlisting}
\begin{block}{The three-step idea}
\begin{enumerate}
\item \textbf{Split:} create groups such as each city.
\item \textbf{Apply:} calculate a sum, mean, count, or other measure.
\item \textbf{Combine:} return a result for each group.
\end{enumerate}
\end{block}
\scenario{A regional manager asks which city generates the most revenue. Group by city and sum revenue, then sort the result.}
\end{frame}

% 31
\begin{frame}[fragile]{GroupBy with several summaries}
\begin{lstlisting}
summary = sales.groupby("city").agg(
    order_count=("order_id", "count"),
    units_sold=("units", "sum"),
    average_revenue=("revenue", "mean"),
    top_revenue=("revenue", "max")
).sort_values("units_sold", ascending=False)
\end{lstlisting}
\begin{itemize}
\item \method{agg()} lets you calculate several statistics in one grouped operation.
\item Named aggregation gives the output columns clear names.
\item Sort the result to answer a decision question.
\end{itemize}
\scenario{A store owner wants order count, total units, average order revenue, and the largest order for each city in one report.}
\end{frame}

% 32
\begin{frame}[fragile]{GroupBy: filter groups or add group context}
\begin{lstlisting}
# Keep cities with at least 3 rows in this example
busy = sales.groupby("city").filter(lambda g: len(g) >= 3)

# Add each city's total revenue to every row in that city
sales["city_total"] = sales.groupby("city")["revenue"].transform("sum")

# Rank orders inside each city
sales["rank_in_city"] = sales.groupby("city")["revenue"].rank(
    method="dense", ascending=False
)
\end{lstlisting}
\begin{itemize}
\item \method{filter}: keep or remove entire groups.
\item \method{transform}: return values aligned to the original rows.
\item \method{rank}: rank values, optionally within each group.
\end{itemize}
\scenario{You want to compare each order with its city's total without shrinking the dataset to one row per city.}
\end{frame}

% 33
\begin{frame}[fragile]{Apply, map, and transform: which one?}
\begin{description}
\item[\method{map}] Best for mapping values in a Series, often using a dictionary.
\item[\method{apply}] Flexible; use when a built-in vectorized operation does not express the task.
\item[\method{transform}] Useful when group calculations must return a result aligned with the original rows.
\item[\method{agg}] Best for summary statistics that reduce each group to a result.
\end{description}
\begin{lstlisting}
sales["city_code"] = sales["city"].map({"Dhaka": "DHK"})
sales.groupby("city")["revenue"].agg(["sum", "mean"])
sales["city_avg"] = sales.groupby("city")["revenue"].transform("mean")
\end{lstlisting}
\pitfall{Do not default to \method{apply()} for everything. Built-in methods are often clearer and faster.}
\end{frame}

% 34
\begin{frame}[fragile]{Combine tables: concatenate}
\begin{lstlisting}
jan = pd.DataFrame({"order_id": [1, 2], "amount": [100, 200]})
feb = pd.DataFrame({"order_id": [3, 4], "amount": [150, 250]})

all_orders = pd.concat([jan, feb], ignore_index=True)
side_by_side = pd.concat([jan, feb], axis=1)
\end{lstlisting}
\begin{itemize}
\item \method{concat(..., axis=0)} stacks tables vertically by default.
\item \method{concat(..., axis=1)} aligns and places columns side by side.
\item \texttt{ignore\_index=True} creates a fresh numbered index when stacking rows.
\end{itemize}
\scenario{The sales team sends one file per month. Stack the monthly files into one annual table before analyzing trends.}
\end{frame}

% 35
\begin{frame}[fragile]{Combine tables: merge like a database JOIN}
\begin{lstlisting}
orders = pd.DataFrame({
    "customer_id": [10, 11, 12],
    "amount": [500, 250, 700]
})
customers = pd.DataFrame({
    "customer_id": [10, 11, 13],
    "name": ["Asha", "Rafi", "Mina"]
})

joined = orders.merge(customers, on="customer_id", how="left")
\end{lstlisting}
\begin{itemize}
\item \texttt{how="inner"}: only matching keys.
\item \texttt{how="left"}: keep all rows from the left table.
\item \texttt{how="right"}: keep all rows from the right table.
\item \texttt{how="outer"}: keep keys from both tables.
\end{itemize}
\scenario{Orders contain customer IDs; a separate customer table contains names. Merge them to create a readable order report.}
\end{frame}

% 36
\begin{frame}[fragile]{Merge safely: keys, suffixes, and validation}
\begin{lstlisting}
result = orders.merge(
    customers,
    on="customer_id",
    how="left",
    validate="many_to_one",
    indicator=True
)
result["_merge"].value_counts()

# If key column names differ:
# left.merge(right, left_on="cust_id", right_on="customer_id")
\end{lstlisting}
\begin{itemize}
\item \texttt{validate} checks the expected relationship, such as many orders to one customer.
\item \texttt{indicator=True} labels rows as left-only, right-only, or matched.
\item Use \texttt{suffixes=("\_orders", "\_customers")} when both tables share non-key column names.
\end{itemize}
\pitfall{A merge can unexpectedly multiply rows if a key appears more than once on both sides. Check key uniqueness and row counts.}
\end{frame}

% 37
\begin{frame}[fragile]{Join and ordered merge}
\begin{lstlisting}
# Join using index labels
left = left.set_index("customer_id")
right = right.set_index("customer_id")
joined = left.join(right, how="left")

# Match each event to the latest earlier timestamp
# quotes.merge_asof(trades, on="time", direction="backward")
\end{lstlisting}
\begin{itemize}
\item \method{join()} is convenient for index-based joins.
\item \method{merge\_asof()} matches nearby ordered values, commonly timestamps.
\item \method{merge\_ordered()} can help combine ordered data such as time-indexed records.
\end{itemize}
\scenario{In market data, match each trade to the most recent quote before that trade. This is an as-of matching problem, not a normal exact-key join.}
\pitfall{\method{merge\_asof()} requires sorted merge keys. Read the documentation carefully before using it in a financial or scientific workflow.}
\end{frame}

% 38
\begin{frame}[fragile]{Reshape: wide versus long data}
\begin{columns}[T,onlytextwidth]
\column{0.49\textwidth}
\begin{block}{Wide}
\begin{tabular}{lrr}
\toprule
Student & Math & Science\\
\midrule
Ava & 80 & 90\\
Rafi & 75 & 85\\
\bottomrule
\end{tabular}
\small Each subject gets a separate column.
\end{block}
\column{0.49\textwidth}
\begin{block}{Long}
\begin{tabular}{llr}
\toprule
Student & Subject & Mark\\
\midrule
Ava & Math & 80\\
Ava & Science & 90\\
Rafi & Math & 75\\
Rafi & Science & 85\\
\bottomrule
\end{tabular}
\small Subject becomes a value in one column.
\end{block}
\end{columns}
\scenario{Long format is often convenient for plotting and grouped analysis. Wide format can be convenient for reports and spreadsheets.}
\end{frame}

% 39
\begin{frame}[fragile]{Reshape with \texttt{melt} and \texttt{pivot}}
\begin{lstlisting}
long = marks.melt(
    id_vars="student",
    var_name="subject",
    value_name="mark"
)

wide = long.pivot(
    index="student",
    columns="subject",
    values="mark"
)
\end{lstlisting}
\begin{itemize}
\item \method{melt()} turns selected columns into rows.
\item \method{pivot()} spreads unique values into columns.
\item \method{pivot()} needs each index/column pair to identify one value.
\end{itemize}
\pitfall{If duplicate pairs exist, \method{pivot()} raises an error. If you need to combine duplicates, use \method{pivot\_table()} with an aggregation function.}
\end{frame}

% 40
\begin{frame}[fragile]{Pivot tables and cross-tabulations}
\begin{lstlisting}
pd.pivot_table(
    sales,
    values="revenue",
    index="city",
    columns="product",
    aggfunc="sum",
    fill_value=0,
    margins=True
)

pd.crosstab(survey["age_band"], survey["answer"])
\end{lstlisting}
\begin{itemize}
\item \method{pivot\_table()} summarizes data across row and column categories.
\item \texttt{aggfunc} selects how repeated values are combined.
\item \method{crosstab()} counts combinations of categories by default.
\end{itemize}
\scenario{A store wants a grid showing total revenue for every city-product combination. A school survey wants counts of answers by age band.}
\end{frame}

% 41
\begin{frame}[fragile]{Other reshaping tools}
\begin{lstlisting}
# Turn each list item into its own row
items = pd.DataFrame({
    "order_id": [1, 2],
    "items": [["pen", "book"], ["bag"]]
})
expanded = items.explode("items")

# Convert categories to indicator columns
encoded = pd.get_dummies(df["payment_type"], prefix="pay")

# Move between hierarchical index levels
# stacked = wide.stack()
# wide_again = stacked.unstack()
\end{lstlisting}
\begin{itemize}
\item \method{explode()} expands list-like values into rows.
\item \method{get\_dummies()} creates 0/1 indicator columns.
\item \method{stack()} and \method{unstack()} move between index levels and columns.
\end{itemize}
\scenario{A single order contains several products in a list. Explode it to analyze product-level demand.}
\end{frame}

% 42
\begin{frame}[fragile]{Dates and times: parse them early}
\begin{lstlisting}
df["order_date"] = pd.to_datetime(df["order_date"], errors="coerce")
df["year"] = df["order_date"].dt.year
df["month"] = df["order_date"].dt.month
df["weekday"] = df["order_date"].dt.day_name()
df["is_weekend"] = df["order_date"].dt.dayofweek >= 5
\end{lstlisting}
\begin{itemize}
\item \method{to\_datetime()} converts date-like values into datetime objects.
\item \method{dt} provides date components and date-related methods for a datetime Series.
\item Parse dates before sorting, filtering, or grouping by time.
\end{itemize}
\scenario{A delivery company wants to compare weekday and weekend order volume.}
\pitfall{Dates can be ambiguous: \texttt{03/04/2026} may mean March 4 or April 3 depending on the source. Confirm the format.}
\end{frame}

% 43
\begin{frame}[fragile]{Time-based filtering and resampling}
\begin{lstlisting}
df = df.sort_values("order_date")
recent = df[df["order_date"] >= "2026-01-01"]

daily = (
    df.set_index("order_date")
      .resample("D")["revenue"]
      .sum()
)
monthly = (
    df.set_index("order_date")
      .resample("MS")["revenue"]
      .sum()
)
\end{lstlisting}
\begin{itemize}
\item \method{resample()} groups time-indexed values into intervals such as days or months.
\item \texttt{D} means daily; \texttt{MS} means month start.
\item Choose the aggregation that fits the metric: sum revenue, mean temperature, or count events.
\end{itemize}
\end{frame}

% 44
\begin{frame}[fragile]{Rolling, expanding, and shifting}
\begin{lstlisting}
df = df.sort_values("date").set_index("date")
df["7_day_avg"] = df["revenue"].rolling("7D").mean()
df["running_total"] = df["revenue"].expanding().sum()
df["previous_day"] = df["revenue"].shift(1)
df["change"] = df["revenue"].diff()
df["pct_change"] = df["revenue"].pct_change()
\end{lstlisting}
\begin{itemize}
\item \method{rolling()}: moving window, such as a 7-day average.
\item \method{expanding()}: calculation over all observations so far.
\item \method{shift()}: align with a previous or next row.
\item \method{diff()} and \method{pct\_change()}: absolute and relative changes.
\end{itemize}
\scenario{A shop uses a rolling average to smooth daily sales fluctuations and a percentage change to compare performance.}
\pitfall{A row-based window of 7 observations is not always seven calendar days. Use a time-based window when gaps in dates matter.}
\end{frame}

% 45
\begin{frame}[fragile]{Time zones and elapsed time}
\begin{lstlisting}
df["timestamp"] = pd.to_datetime(df["timestamp"], utc=True)
df["local_time"] = df["timestamp"].dt.tz_convert("Asia/Dhaka")

start = pd.to_datetime("2026-01-01")
end = pd.to_datetime("2026-01-08")
elapsed = end - start
\end{lstlisting}
\begin{itemize}
\item Time zone handling matters when data crosses regions or daylight-saving changes.
\item UTC is a useful common reference for storing event timestamps.
\item A time difference is a \texttt{Timedelta}.
\end{itemize}
\scenario{An international support team compares tickets from different time zones. Normalize timestamps before measuring response times.}
\pitfall{Do not localize a timestamp and convert a timestamp interchangeably. Localization assigns a time zone to a naive clock reading; conversion changes the displayed time zone of an aware timestamp.}
\end{frame}

% 46
\begin{frame}[fragile]{Categorical data: when repeated labels dominate}
\begin{lstlisting}
df["priority"] = df["priority"].astype(
    pd.CategoricalDtype(
        categories=["Low", "Medium", "High"],
        ordered=True
    )
)
df.sort_values("priority")
\end{lstlisting}
\begin{itemize}
\item A categorical dtype represents values from a known set of labels.
\item Ordered categories can express a meaningful order such as Low, Medium, High.
\item Categories may reduce memory use when many values repeat.
\end{itemize}
\scenario{A ticketing system needs priorities sorted by business meaning, not alphabetically.}
\pitfall{Define categories carefully: unexpected labels may become missing during conversion.}
\end{frame}

% 47
\begin{frame}[fragile]{Plot directly from pandas}
\begin{lstlisting}
import matplotlib.pyplot as plt

sales.groupby("city")["revenue"].sum().plot(
    kind="bar", title="Revenue by city", ylabel="Revenue"
)
plt.tight_layout()
plt.show()

sales["revenue"].plot(kind="hist", bins=10, title="Order revenue")
plt.show()
\end{lstlisting}
\begin{itemize}
\item Common kinds: \texttt{line}, \texttt{bar}, \texttt{barh}, \texttt{hist}, \texttt{box}, \texttt{scatter}.
\item Line charts are useful for ordered time trends; bars compare categories.
\item Histograms show distributions; box plots help spot spread and possible outliers.
\end{itemize}
\scenario{Show a manager which city contributes the most revenue, or whether most orders are small with a few very large ones.}
\end{frame}

% 48
\begin{frame}[fragile]{Choose a chart that answers the question}
\begin{tabular}{p{0.24\textwidth}p{0.32\textwidth}p{0.34\textwidth}}
\toprule
\textbf{Question} & \textbf{Useful chart} & \textbf{Example}\\
\midrule
How does it change over time? & Line & Daily sales over a month\\
Which category is larger? & Bar & Revenue by city\\
How are values distributed? & Histogram & Delivery times\\
Are there unusual values? & Box plot & Product prices\\
Do two numeric values move together? & Scatter & Ad spend vs sales\\
\bottomrule
\end{tabular}
\begin{block}{A good chart}
Use a descriptive title, readable labels, appropriate units, and a scale that does not mislead the viewer.
\end{block}
\end{frame}

% 49
\begin{frame}[fragile]{Style a table for a report}
\begin{lstlisting}
summary = sales.groupby("city", as_index=False)["revenue"].sum()

styled = (
    summary.style
    .format({"revenue": "{:,.2f}"})
    .highlight_max(subset=["revenue"])
    .set_caption("Revenue by city")
)
\end{lstlisting}
\begin{itemize}
\item \method{DataFrame.style} creates a display-oriented Styler object in notebook or HTML workflows.
\item \method{format()} controls how values appear.
\item Highlighting can draw attention to important values.
\end{itemize}
\pitfall{Styling is for presentation. It does not change the underlying numeric values, and support differs by output format.}
\end{frame}

% 50
\begin{frame}[fragile]{Useful comparisons and aligned arithmetic}
\begin{lstlisting}
left = pd.Series([10, 20], index=["A", "B"])
right = pd.Series([3, 7], index=["B", "C"])

left + right
left.add(right, fill_value=0)
\end{lstlisting}
\begin{itemize}
\item pandas aligns Series and DataFrame values by labels during many operations.
\item When labels do not match, the result may contain missing values.
\item Methods such as \method{add()}, \method{sub()}, \method{mul()}, and \method{div()} can accept \texttt{fill\_value}.
\end{itemize}
\scenario{Compare two reports whose product rows are in different orders. Labels help pandas match the correct products rather than blindly matching row positions.}
\end{frame}

% 51
\begin{frame}[fragile]{Compare datasets and find changes}
\begin{lstlisting}
old = pd.DataFrame({"price": [100, 200]}, index=["A", "B"])
new = pd.DataFrame({"price": [110, 200]}, index=["A", "B"])

old.compare(new)
old.equals(new)
\end{lstlisting}
\begin{itemize}
\item \method{compare()} highlights differences between similarly shaped objects.
\item \method{equals()} returns a boolean for exact equality under pandas' comparison rules.
\item Align rows and columns first if the datasets have different structures.
\end{itemize}
\scenario{A team receives a revised inventory file and needs to find which product prices changed before publishing it.}
\end{frame}

% 52
\begin{frame}[fragile]{Indexes: labels are powerful, but not magic}
\begin{lstlisting}
df.index.is_unique
df.columns.is_unique
df = df.set_index("customer_id", verify_integrity=True)
df = df.reset_index()
\end{lstlisting}
\begin{itemize}
\item An index labels rows and can support selection, alignment, and joins.
\item Index labels do not have to be consecutive integers.
\item Duplicate index or column labels can make some operations confusing.
\item \texttt{verify\_integrity=True} can catch duplicate keys when setting an index.
\end{itemize}
\scenario{A customer ID should identify one customer in a customer master table. Verify uniqueness before using it as a lookup index.}
\end{frame}

% 53
\begin{frame}[fragile]{MultiIndex: more than one label level}
\begin{lstlisting}
grouped = sales.groupby(["city", "product"])["revenue"].sum()
grouped.loc[("Dhaka", "Pen")]

# Convert index levels back to columns
flat = grouped.reset_index()
\end{lstlisting}
\begin{itemize}
\item A MultiIndex stores multiple levels of labels on an axis.
\item It is useful for grouped results, panel-like data, or nested categories.
\item \method{reset\_index()} makes the levels easier to export or join as ordinary columns.
\end{itemize}
\scenario{You want revenue for each city-product pair, not just one summary per city.}
\end{frame}

% 54
\begin{frame}[fragile]{Performance: make pandas work smarter}
\begin{itemize}
\item Load only the columns and rows needed for the task.
\item Prefer vectorized operations and built-in aggregations over row-by-row Python loops.
\item Use suitable dtypes; repeated labels may benefit from categorical storage.
\item Avoid repeatedly concatenating one row at a time. Collect pieces, then concatenate once.
\item Check memory use with \method{memory\_usage(deep=True)}.
\item Measure before optimizing; the slowest-looking line is not always the bottleneck.
\end{itemize}
\begin{lstlisting}
parts = []
for file in files:
    parts.append(pd.read_csv(file, usecols=["date", "sales"]))
combined = pd.concat(parts, ignore_index=True)
\end{lstlisting}
\end{frame}

% 55
\begin{frame}[fragile]{Work with data too large for memory}
\begin{lstlisting}
totals = []
for chunk in pd.read_csv("huge_orders.csv", chunksize=100_000):
    chunk["revenue"] = chunk["units"] * chunk["price"]
    totals.append(chunk.groupby("city")["revenue"].sum())

final = pd.concat(totals).groupby(level=0).sum()
\end{lstlisting}
\begin{itemize}
\item \texttt{chunksize} reads a file in manageable pieces.
\item Summarize each chunk, then combine the partial summaries.
\item For workflows requiring large joins or complex processing beyond memory, consider a database or a scalable dataframe tool.
\end{itemize}
\pitfall{Chunking only works directly when you can process parts independently or combine partial results correctly. Medians and many complex operations need more careful strategies.}
\end{frame}

% 56
\begin{frame}[fragile]{Copying, mutation, and chained assignment}
\begin{lstlisting}
# Explicitly create a separate working DataFrame
dhaka_sales = sales.loc[sales["city"] == "Dhaka"].copy()
dhaka_sales["flag"] = True

# Set selected cells in one operation
sales.loc[sales["revenue"] > 500, "priority"] = "High"
\end{lstlisting}
\begin{itemize}
\item Make the intended target of an assignment clear.
\item Avoid chained patterns such as \texttt{df[mask]["col"] = value}.
\item Modern pandas uses Copy-on-Write behavior; code should not rely on accidental mutation of another object.
\end{itemize}
\pitfall{If a transformation must be independent, use \method{copy()} explicitly and keep assignments direct. Consult the Copy-on-Write guide for version-specific behavior.}
\end{frame}

% 57
\begin{frame}[fragile]{Options and display settings}
\begin{lstlisting}
pd.set_option("display.max_columns", 20)
pd.set_option("display.max_rows", 30)
pd.set_option("display.float_format", "{:.2f}".format)

# Temporarily change settings
with pd.option_context("display.max_columns", 10):
    print(df)
\end{lstlisting}
\begin{itemize}
\item Display options change how objects are shown, not the actual stored data.
\item \method{option\_context()} restores previous options when the block ends.
\item Keep display settings in notebooks or scripts where they improve readability.
\end{itemize}
\end{frame}

% 58
\begin{frame}[fragile]{Common mistakes and their fixes}
\begin{tabular}{p{0.42\textwidth}p{0.47\textwidth}}
\toprule
\textbf{Mistake} & \textbf{Better approach}\\
\midrule
Comparing text numbers to numeric thresholds & Convert types and inspect failed conversions\\
Using \texttt{and}/\texttt{or} between Series & Use \texttt{\&}/\texttt{|} with parentheses\\
Assuming row position equals index label & Choose \texttt{loc} or \texttt{iloc} deliberately\\
Filling every missing value with zero & Use domain knowledge to choose a strategy\\
Merging without checking duplicate keys & Check uniqueness and use \texttt{validate}\\
Looping through every row for simple math & Prefer vectorized column operations\\
Deleting every repeated row automatically & Define what makes a record truly duplicate\\
\bottomrule
\end{tabular}
\end{frame}

% 59
\begin{frame}[fragile]{A complete cleaning workflow}
\begin{lstlisting}
df = pd.read_csv("raw_orders.csv")

# 1. Understand the data
print(df.shape)
df.info()
print(df.isna().sum())

# 2. Standardize labels and types
df.columns = df.columns.str.strip().str.lower()
df["order_date"] = pd.to_datetime(df["order_date"], errors="coerce")
df["price"] = pd.to_numeric(df["price"], errors="coerce")

# 3. Resolve known quality issues
df = df.drop_duplicates(subset=["order_id"])
df = df.dropna(subset=["order_id", "price"])

# 4. Add a useful measure and save a new file
df["revenue"] = df["units"] * df["price"]
df.to_csv("clean_orders.csv", index=False)
\end{lstlisting}
\small This is a template, not a universal recipe: the right missing-value and duplicate rules depend on the meaning of the data.
\end{frame}

% 60
\begin{frame}[fragile]{A complete analysis workflow}
\begin{lstlisting}
df = pd.read_csv("clean_orders.csv")
df["order_date"] = pd.to_datetime(df["order_date"])

city_report = (
    df.groupby("city", as_index=False)
      .agg(
          orders=("order_id", "nunique"),
          units=("units", "sum"),
          revenue=("revenue", "sum")
      )
      .sort_values("revenue", ascending=False)
)

monthly = (
    df.set_index("order_date")
      .resample("MS")["revenue"]
      .sum()
)

city_report.to_csv("city_report.csv", index=False)
\end{lstlisting}
\scenario{This workflow turns raw order data into two useful outputs: a city performance report and a monthly revenue series.}
\end{frame}

% 61
\begin{frame}[fragile]{Mini project: student performance analysis}
\begin{block}{Scenario}
A school shares a CSV with student ID, class, subject, mark, attendance, and exam date. The head teacher wants a trustworthy summary.
\end{block}
\begin{enumerate}
\item Inspect dimensions, types, duplicates, and missing values.
\item Convert marks, attendance, and exam dates to correct types.
\item Check impossible marks and attendance values.
\item Calculate class average, median, and pass rate.
\item Compare subjects and plot marks distributions.
\item Export a summary and document the cleaning decisions.
\end{enumerate}
\begin{block}{Practice prompts}
Which class has the highest average? Which subject has the widest spread? Are absent students represented as missing, zero, or a separate status?
\end{block}
\end{frame}

% 62
\begin{frame}[fragile]{Mini project: public health or environment data}
\begin{block}{Scenario}
A dataset contains district, date, reported cases, rainfall, temperature, and source.
\end{block}
\begin{itemize}
\item Parse dates and standardize district names.
\item Identify missing periods and duplicate reports.
\item Group by district and month.
\item Use rolling averages to smooth daily fluctuations.
\item Merge rainfall data by a well-defined key and date.
\item Plot trends and clearly label gaps and uncertainty.
\end{itemize}
\pitfall{A pattern or correlation does not prove that one variable caused another. Protect personal data and avoid publishing identifying records.}
\end{frame}

% 63
\begin{frame}[fragile]{Mini project: e-commerce investigation}
\begin{block}{Scenario}
An online shop wants to learn which products and cities contribute to revenue and where customers abandon deliveries.
\end{block}
\begin{enumerate}
\item Clean order status, product names, prices, and timestamps.
\item Remove only confirmed duplicate records.
\item Calculate revenue using the business's documented rule.
\item Merge order lines with product and customer tables.
\item Group by city, product, month, and delivery status.
\item Create a pivot table and two or three decision-focused charts.
\end{enumerate}
\begin{block}{Deliverable}
A reproducible notebook, cleaned dataset, summary CSV, charts, and a short explanation of limitations.
\end{block}
\end{frame}

% 64
\begin{frame}[fragile]{Practice challenge: choose the right method}
\begin{enumerate}
\item Show only rows where price exceeds 500.
\item Find the five largest orders by revenue.
\item Count missing values in each column.
\item Combine January and February order files.
\item Add customer names to order records using customer ID.
\item Calculate total revenue for each city.
\item Convert daily revenue to monthly revenue.
\item Turn subject columns into rows for plotting.
\item Find rows duplicated by email, but keep all copies visible for review.
\item Add each city's total revenue back to every order.
\end{enumerate}
\small Suggested methods: boolean filtering, \texttt{nlargest}, \texttt{isna().sum()}, \texttt{concat}, \texttt{merge}, \texttt{groupby}, \texttt{resample}, \texttt{melt}, \texttt{duplicated(keep=False)}, \texttt{transform}.
\end{frame}

% 65
\begin{frame}[fragile]{A practical method-selection guide}
\begin{tabular}{p{0.32\textwidth}p{0.57\textwidth}}
\toprule
\textbf{Need} & \textbf{Start with}\\
\midrule
Load a file & \texttt{read\_csv}, \texttt{read\_excel}, \texttt{read\_parquet}\\
Inspect quality & \texttt{head}, \texttt{info}, \texttt{describe}, \texttt{isna}\\
Choose rows & Boolean masks, \texttt{loc}, \texttt{query}\\
Clean text & \texttt{str.strip}, \texttt{str.lower}, \texttt{str.replace}\\
Summarize groups & \texttt{groupby} + \texttt{agg}\\
Combine files & \texttt{concat}\\
Match tables by keys & \texttt{merge}\\
Change table shape & \texttt{melt}, \texttt{pivot}, \texttt{pivot\_table}\\
Analyze time & \texttt{to\_datetime}, \texttt{dt}, \texttt{resample}, \texttt{rolling}\\
Reduce memory/time & select columns, efficient dtypes, chunking, measure first\\
\bottomrule
\end{tabular}
\end{frame}

% 66
\begin{frame}[fragile]{Official documentation: the source of truth}
\begin{itemize}
\item \href{https://pandas.pydata.org/docs/}{\textbf{pandas documentation home}} — current documentation and release links.
\item \href{https://pandas.pydata.org/docs/getting_started/}{\textbf{Getting started}} — installation, beginner tutorials, and common tasks.
\item \href{https://pandas.pydata.org/docs/user_guide/}{\textbf{User Guide}} — detailed explanations of core topics.
\item \href{https://pandas.pydata.org/docs/user_guide/10min.html}{\textbf{10 minutes to pandas}} — a compact overview of fundamentals.
\item \href{https://pandas.pydata.org/docs/reference/index.html}{\textbf{API reference}} — exact method parameters and return behavior.
\item \href{https://pandas.pydata.org/docs/user_guide/merging.html}{\textbf{Merge, join, concatenate}} — combining tables.
\item \href{https://pandas.pydata.org/docs/user_guide/reshaping.html}{\textbf{Reshaping and pivot tables}} — changing table layout.
\item \href{https://pandas.pydata.org/docs/user_guide/timeseries.html}{\textbf{Time series}} — dates, offsets, resampling, and time zones.
\end{itemize}
\small pandas changes over time. Check the documentation for the version installed in your environment, especially for dtype and Copy-on-Write behavior.
\end{frame}

% 67
\begin{frame}[fragile]{What does “beginner to advanced” mean?}
\begin{columns}[T,onlytextwidth]
\column{0.32\textwidth}
\begin{block}{Beginner}
\begin{itemize}
\item Create/read tables
\item Inspect columns
\item Select and filter
\item Calculate columns
\item Handle missing values
\end{itemize}
\end{block}
\column{0.32\textwidth}
\begin{block}{Intermediate}
\begin{itemize}
\item Group and aggregate
\item Clean text and dates
\item Merge datasets
\item Reshape tables
\item Visualize findings
\end{itemize}
\end{block}
\column{0.32\textwidth}
\begin{block}{Advanced practice}
\begin{itemize}
\item Time-series windows
\item MultiIndex
\item Memory and chunking
\item Reliable merge checks
\item Reproducible workflows
\end{itemize}
\end{block}
\end{columns}
\begin{center}
\Large\textcolor{Navy}{The goal is not to memorize every method.}\\
\large\textcolor{Coral}{The goal is to choose the right method for the real problem.}
\end{center}
\end{frame}

% 68
\begin{frame}[plain,fragile]
\begin{center}
\vspace{1.3cm}
{\Huge\bfseries\textcolor{Navy}{Your data journey starts here}}\\[0.5cm]
{\Large\textcolor{Blue}{Inspect \(\rightarrow\) Clean \(\rightarrow\) Transform \(\rightarrow\) Analyze \(\rightarrow\) Explain}}\\[1cm]
\begin{block}{Final advice}
Practice with real datasets, keep your raw data unchanged, write down assumptions, verify your results, and use the official pandas documentation whenever you need details.
\end{block}
\vspace{0.4cm}
\href{https://pandas.pydata.org/docs/}{pandas.pydata.org/docs}
\end{center}
\end{frame}

\end{document}
